Formem un circuit amb un cable metàl·lic en forma de U i tanquem la part superior amb una barra metàl·lica que pot lliscar lliurement, tal com mostra la figura. L'amplària de la U és de 0,25~m, i inicialment l'alçària del circuit és de 0,30~m. A continuació, desplacem la barra cap avall a una velocitat constant de $5,00\ \mathrm{m\,s^{-1}}$. Aquest circuit es troba situat en una regió de l'espai on hi ha un camp magnètic de mòdul $B = 2,00\ \mathrm{T}$, perpendicular al pla que forma el circuit i amb sentit cap enfora del paper.

\figura[0.35]{fig1.png}

\begin{apartat}{1,25}
Calculeu l'expressió del flux magnètic del circuit en funció del temps i la força electromotriu induïda al circuit. Indiqueu el sentit de circulació del corrent induït i justifiqueu la resposta.
\end{apartat}

\begin{apartat}{1,25}
Suposem que la barra presenta una resistència de $50,0\ \Omega$ i la resta del circuit no presenta cap resistència. Calculeu la intensitat elèctrica que circularà pel circuit. Calculeu la força magnètica que actua sobre la barra i representeu-la en la figura.
\end{apartat}
